\begin{lemma}\label{lem:decic91-relative-formal-fibre}
Let $R=\mathbb Z_2$, and let $Y/R$ be a smooth proper modular curve of odd level.
Let $\tau$ be a diamond automorphism of order $3$, put $Z=Y/\langle\tau\rangle$,
and write $\pi:Y\longrightarrow Z$ for the quotient map.
Assume that $Z/R$ is smooth.
Suppose $(\tau-1)J_Y(\mathbb Q)$ is finite.
Let $A$ be an effective rational divisor of degree $d$ on the generic fibre.
Let $C$ be an effective divisor over $R$ with $\overline A=\overline C$, and suppose
$(\tau-1)[C-dc_0]$ is torsion in $J_Y(\mathbb Q_2)$, where $c_0$ is a rational cusp.
Write
\[
 C=C_0+\pi^*E,\qquad \deg E=k,
\]
and suppose the map $E'\mapsto C_0+\pi^*E'$ is a closed formal immersion
at $\overline E$. If the restriction of
\[
 U=(\operatorname{Cartier}-1)(\tau^*-1)H^0(Y_{\mathbb F_2},\Omega^1)
\]
to $\overline C$ has rank $d-k$, then $A$ contains $C_0$ as a divisor over $\mathbb Q_2$.
\end{lemma}
\begin{proof}
Put $\alpha=\tau-1$ and $M=T_2-\langle2\rangle$ on the full Jacobian $J_Y$.
The class $\alpha[A-C]$ is torsion and has zero reduction.
Lemma~\ref{lem:dyadic-formal-Hecke-kernel} gives $M\alpha[A-C]=0$.
Thus $A$ belongs to the formal fibre through $C$ of
\[
 \Psi:Y^{(d)}\longrightarrow J_Y,\qquad
 D\longmapsto M\alpha[D-dc_0].
\]
By Eichler--Shimura, the cotangent image of $\Psi$ is $U$;
the invertible diamond factor does not change this image.

Choose $d-k$ local coordinates of $J_Y$ whose pullbacks have independent
differentials at $\overline C$.
Their common fibre is formally smooth over $R$ of relative dimension $k$.
The formal family $C_0+\pi^*Z^{(k)}$ lies in the full fibre of $\Psi$,
since $(\tau-1)\pi^*=0$.
Its tangent map to the chosen smooth fibre is an isomorphism between
$k$-dimensional spaces. It is therefore a formal isomorphism.
The full fibre equals this family, proving the assertion.
\end{proof}

We use the lemma only in the following two situations, where the stated
closed immersion can be checked directly. At an unramified orbit of three
cusps, subtract the minimum of their multiplicities. At least one residual
multiplicity is zero, and the divisor at that point recovers the monic
polynomial defining the moving divisor downstairs. The same argument applies
to a cubic cusp whose three geometric points form one orbit.
At a tame fixed quadratic point $P$, choose, after unramified base change,
a parameter $z$ with $\tau z=\zeta_3z$. The quotient parameter is $z^3$.
Near $P+\pi^*(\pi(P))=4P$, the family is given by
$z(z^3-a)$, whose coefficient of $z$ recovers $a$.
These descriptions give closed formal immersions over the split unramified
extension, and hence over $R$ by descent.
